\section{Symplectic components at two}\label{sec:symplectic}

For a symplectic component, signs on a decomposition into planes leave
room for a quaternion four-group.  Its two-dimensional space of local
chain coefficients provides the extra factor needed in the smallest
case.  We construct the cycle directly inside the selected faithful
component extension, then apply Lemma~\ref{lem:propagation} at a
supported full flag.

\begin{proposition}\label{prop:symplectic}
A minimal-order counterexample to \eqref{eq:HQC} at $p=2$ cannot have a
simple component isomorphic to $\PSp_{2n}(q)$, where $n\geq3$ and $q$
has characteristic at least five.
\end{proposition}

Throughout this section, $V$ is a $2n$-dimensional vector space over
$\F_q$ with nondegenerate alternating form $\omega$, and
$L=\PSp(V)$.  Write $M=\PGSp(V)$ for the projective group of linear
symplectic similitudes.  In this range the automorphism theorem
\citep[Section~2 and statements~3.2--3.6]{Steinberg1960} gives
$\Aut(L)=M\rtimes\Gal(\F_q/\F_\ell)$, where $\ell$ is the
characteristic; there is no additional diagram automorphism.  Thus
$\Aut(L)/M$ is cyclic.  We use projective classes of semilinear
similitudes for the elements of this group, as in
Section~\ref{sec:reductions}.

\subsection{Plane signs and the quaternion action}

Let $\mathcal D=\{V_1,\ldots,V_n\}$ be an unordered orthogonal
decomposition of $V$ into nondegenerate planes.  Its linear sign group
and projective sign group are
\[
\widehat S(\mathcal D)
 =\{\diag(\epsilon_1\id_{V_1},\ldots,
                   \epsilon_n\id_{V_n}):\epsilon_i\in\{1,-1\}\},
\qquad
S(\mathcal D)=\widehat S(\mathcal D)/\langle-\id_V\rangle.
\]
All these signs are symplectic, and $S(\mathcal D)\leq L$ has rank
$n-1$.

A \emph{quaternion four-group compatible with $\mathcal D$} is a
subgroup $T\leq L$ generated by the projective classes of two
operators $I,J\in\Sp(V)$ such that
\begin{equation}\label{eq:symp-quaternion}
I^2=J^2=-\id_V,\qquad IJ=-JI,
\end{equation}
and both operators preserve every $V_i$.  On each plane these
operators generate its full matrix algebra.  Indeed, after extending
scalars to an algebraic closure, $I$ has two distinct eigenvalues and
$J$ interchanges the two eigenlines.  The resulting diagonal and
off-diagonal matrix units span the full matrix algebra; its dimension
is therefore already four over $\F_q$.

Such pairs exist over every finite field of odd order.  On a plane
choose
\[
I=\begin{pmatrix}0&1\\-1&0\end{pmatrix},\qquad
J=\begin{pmatrix}a&b\\b&-a\end{pmatrix},
\qquad a^2+b^2=-1.
\]
The equation has a solution: if $-1$ is a square this follows by
factoring the left side over the field, and otherwise it is the
surjectivity of the norm from its quadratic extension.  These matrices
have determinant one and satisfy \eqref{eq:symp-quaternion}.
Repeating the construction on orthogonal planes gives the required
operators on $V$.  Each of $I,J,IJ$ is nonscalar on every plane, so
$T\cap S(\mathcal D)=1$.  Consequently
\begin{equation}\label{eq:symp-seed}
A_0(\mathcal D,T)=S(\mathcal D)\times T\leq L,
\qquad \rk A_0(\mathcal D,T)=n+1.
\end{equation}

\begin{lemma}\label{lem:symp-centralizer}
For $n\geq3$ and a pair $(\mathcal D,T)$ as above,
\[
C_M\bigl(A_0(\mathcal D,T)\bigr)=A_0(\mathcal D,T).
\]
Every semilinear projective transformation centralizing
$S(\mathcal D)$ preserves each plane of $\mathcal D$ individually.
\end{lemma}

\begin{proof}
The weights of $\widehat S(\mathcal D)\cong\F_2^n$ on $V$ are its
distinct coordinate characters $e_1,\ldots,e_n$, with weight spaces
$V_1,\ldots,V_n$.  A projective centralizer, represented by $g$, satisfies
\[
g d g^{-1}=c(d)d\qquad(d\in\widehat S(\mathcal D)),
\]
where $c$ is a character with values in $\{1,-1\}$: these values follow
by squaring, and multiplicativity follows by taking products.  This
holds also for semilinear $g$, since field automorphisms fix the signs.
The induced permutation of weight spaces must therefore preserve the
set $\{e_1,\ldots,e_n\}$ by translation by a single character $t$.
If $t\ne0$ and $e_i+t=e_j$, then $i\ne j$ and $t=e_i+e_j$.
For a third index $k$, the character $e_k+t$ has three nonzero
coordinates and is not a coordinate character.  This contradiction
proves individual preservation of all the planes.

Now let $g$ be linear and centralize $A_0$ projectively.  Projective
commutation with $I$ and $J$ gives
\[
gIg^{-1}=\epsilon I,\qquad gJg^{-1}=\eta J,
\qquad \epsilon,\eta\in\{1,-1\},
\]
because $I^2=J^2=-\id_V$.  Conjugation by $I$ changes the sign of $J$,
and conjugation by $J$ changes the sign of $I$.  Multiplying $g$ by
one of $1,I,J,IJ$ therefore makes it commute exactly with both
operators.  It still preserves every plane.  The full matrix-algebra
calculation above now shows that its restriction to $V_i$ is
$\lambda_i\id_{V_i}$.  The similitude multiplier is $\lambda_i^2$
on each nondegenerate plane, so all $\lambda_i^2$ are equal.
After multiplication by the common scalar $\lambda_1^{-1}$, these
restrictions are independent signs.  Hence the projective class of
the original $g$ lies in $S(\mathcal D)T$.  The reverse inclusion
holds because $A_0$ is abelian.
\end{proof}

For fixed $T$, the compatible decompositions have the form required
by the symmetric case of Theorem~\ref{thm:frame-bound}.  Here are the
details of that identification.  The algebra generated by $I$ and $J$
is $\End(U)\cong M_2(\F_q)$, where $U$ is its two-dimensional simple
module.  Complete reducibility, or the matrix units of this algebra,
gives
\begin{equation}\label{eq:symp-multiplicity}
V\cong U\otimes_{\F_q} W,\qquad \dim W=n.
\end{equation}
Choose a nonzero invariant alternating form $\omega_U$ on $U$.
Invariant bilinear forms on $U$ form a one-dimensional space: under
the identification $U\cong U^*$ supplied by $\omega_U$, a second
invariant form corresponds to an operator commuting with $I,J$,
and hence to a scalar.  It follows that the original form is
\[
\omega=\omega_U\otimes B_W
\]
for a bilinear form $B_W$ on $W$.  Alternation of $\omega$ makes
$B_W$ symmetric, and nondegeneracy of $\omega$ makes $B_W$
nondegenerate.  Submodules of $U\otimes W$ are exactly the spaces
$U\otimes W'$ with $W'\leq W$.  In particular, the nondegenerate
invariant planes are $U\otimes w$, where $w$ is a nondegenerate line
of $W$, and two such planes are orthogonal exactly when their lines
are orthogonal for $B_W$.  Thus compatible plane decompositions are
in bijection with all orthogonal nondegenerate line frames of the
symmetric space $(W,B_W)$.  Both line norm types are allowed.

\subsection{Maximal configurations and field descent}

Suppose, toward Proposition~\ref{prop:symplectic}, that $G$ is a
minimal counterexample with the component $L$.  Consider all elementary
abelian subgroups $A\leq N_G(L)$ such that
\begin{equation}\label{eq:symp-configurations}
A\text{ contains some }A_0(\mathcal D,T),\qquad
LA\text{ acts faithfully on }L,\qquad O_2(C_G(LA))=1.
\end{equation}
This family is closed under elementary overgroups that are faithful
configurations, since such an overgroup still contains the same seed.
It is nonempty: take $A=A_0\leq L$ and apply
Lemma~\ref{lem:component-core} to obtain $O_2(C_G(L))=1$.
Choose a member $A$ of maximum rank, put $m=\rk A$, and write
$H=LA$.  Identify this faithful group with its image in $\Aut(L)$.

Since $A$ centralizes its subgroup $A_0$, Lemma~\ref{lem:symp-centralizer}
gives $A\cap M=A_0$.  The quotient $\Aut(L)/M$ is cyclic, so its
elementary abelian subgroups have rank at most one.  Therefore
\begin{equation}\label{eq:symp-ranks}
m=n+1\quad\text{or}\quad m=n+2.
\end{equation}
In the first case $A=A_0$ and $H=L$.  In the second,
\[
A=A_0\times\langle b\rangle,
\qquad H/L\cong C_2,\qquad H\cap M=L,
\]
where $b$ induces the involution $\sigma$ of $\F_q/\F_s$ and
$q=s^2$.  The field parameter $s$ is at least five.

The next normalization is essential for the construction.  We will
require the field axis to commute \emph{exactly} with the quaternion
lifts, not merely with their projective classes.

\begin{lemma}\label{lem:symp-descent}
Let $T=\langle\overline I,\overline J\rangle$ satisfy
\eqref{eq:symp-quaternion}, and let the projective semilinear
involution $b$ centralize $S(\mathcal D)T$ and induce
$\sigma\in\Gal(\F_q/\F_s)$, where $q=s^2$.  Replacing $b$ by
$tb$ for some $t\in T$ leaves $\langle S(\mathcal D),T,b\rangle$
unchanged and permits a semilinear representative $B$ with
\[
B^2=\id_V,\qquad BI=IB,\qquad BJ=JB.
\]
After scaling $\omega$, its restriction to $V_0=\{v:Bv=v\}$ is a
nondegenerate alternating form over $\F_s$, and $I,J$ descend to
symplectic operators on $V_0$.  For fixed $T,b$ satisfying these exact
commutation conditions, the compatible plane decompositions are all
orthogonal nondegenerate line frames of a symmetric
$n$-dimensional space over $\F_s$.
\end{lemma}

\begin{proof}
For any semilinear representative of $b$, conjugation of $I$ and $J$
differs from the identity by two signs, since their squares are
$-\id_V$ and $\sigma$ fixes $-1$.  As in
Lemma~\ref{lem:symp-centralizer}, multiplication by one of
$1,I,J,IJ$ removes both signs.  Its projective class is still an
involution: the classes $b$ and $t$ commute and have order two.
Denote the resulting representative temporarily by $B$.  Since
$B^2=c\id_V$, the identity $B B^2=B^2 B$ implies $\sigma(c)=c$.
The norm $\F_q^\times\to\F_s^\times$ is surjective, so choose
$a$ with $a\sigma(a)=c^{-1}$.  Replacing $B$ by $aB$ makes its
square the identity and preserves exact commutation with $I,J$.

Write the semilinear similitude identity as
\[
\omega(Bv,Bw)=\mu\,\sigma(\omega(v,w)).
\]
The equality $B^2=\id_V$ gives $\mu\sigma(\mu)=1$.  Choose
$t\in\F_q^\times$ with $\sigma(t)/t=\mu$, by the finite-field
form of Hilbert's Theorem~90
\citep[Corollary~5.25]{MilneFields2022}.  Then $\omega'=t\omega$ satisfies
\[
\omega'(Bv,Bw)=\sigma(\omega'(v,w)).
\]
The fixed-vector descent statement in
Section~\ref{subsec:linear-compatible} identifies $V$ with
$\F_q\otimes_{\F_s}V_0$.  The form $\omega'$ takes values in
$\F_s$ on $V_0$, and its restriction is nondegenerate because its
scalar extension is $\omega'$.  Exact commutation makes $I,J$
operators on $V_0$ preserving this form and satisfying the same
quaternion equations.

Each original plane is $B$-stable by
Lemma~\ref{lem:symp-centralizer}, and hence descends to a
nondegenerate symplectic plane over $\F_s$.  Applying
\eqref{eq:symp-multiplicity} over $\F_s$ gives
$V_0\cong U_0\otimes W_0$ with an alternating form on $U_0$ and a
nondegenerate symmetric form on $W_0$.  The preceding submodule
argument identifies all the compatible decompositions with all the
orthogonal line frames of $W_0$.  Conversely, every such frame
extends to a decomposition compatible with $T$ and $B$.  This proves
the assertion in both directions.
\end{proof}

We now specify the finite family of data from which the cycle will be
formed.  If $m=n+1$, a datum is any pair $(\mathcal D,T)$ as in
\eqref{eq:symp-seed}; its full elementary group is $E=S(\mathcal D)T$.
If $m=n+2$, a datum is a triple $(\mathcal D,T,b)$, where
$(\mathcal D,T)$ has that meaning, $b\in H\setminus L$ is an
involution preserving each plane, and a semilinear representative of
$b$ commutes exactly with the lifts $I,J$.  The full elementary group
is then
\begin{equation}\label{eq:symp-full-data}
E=S(\mathcal D)\times T\times\langle b\rangle.
\end{equation}
The exact condition is independent of the scalar representative of
$b$ and of the choices of the signs of $I,J$.  It also holds for
every element of their quaternion group once it holds for $I,J$.
Lemma~\ref{lem:symp-descent} shows that this family is nonempty, since
the original field axis may be shifted inside $T\leq L$.

For each fixed decoration $T$, or fixed pair $(T,b)$, we use every
compatible frame.  The preceding arguments identify these frames
with the symmetric spaces in Theorem~\ref{thm:frame-bound}, over
\begin{equation}\label{eq:symp-field-parameter}
u=q\quad\text{if }m=n+1,\qquad
u=s\quad\text{if }m=n+2.
\end{equation}
We impose no restriction on the isometry type of the multiplicity
space as the decorations vary.  The frame bound is uniform in that
type; this matters because changing a field axis can change the
descended symmetric form.  Coarsening two lines of the multiplicity
space coarsens their two invariant planes, and the equal-sign
subgroups are unchanged by the identification.  Thus the face
identifications required by the frame theorem hold in the actual
elementary subgroup poset.

\subsection{Completion counts}

To apply Lemma~\ref{lem:quaternion-decoration}, we need lower bounds
for two kinds of completions: restoring a quaternion four-group
from one of its lines, and restoring a field axis.  We keep all the
remaining data fixed in each count.

\begin{lemma}\label{lem:symp-completions}
Let $u$ be as in \eqref{eq:symp-field-parameter}, and let
$\chi_u$ denote the quadratic character of $\F_u^\times$.
A fixed plane frame and retained line of a compatible quaternion
four-group, together with a fixed admissible field axis when present,
have at least
\[
k_T=\frac{(u-\chi_u(-1))^n}{4}\geq4^{n-1}
\]
completions to admissible full data.  In the field case, every
retained pair $(\mathcal D,T)$ that occurs has at least
$k_b=2^{n-1}$ admissible field-axis completions in $H\setminus L$.
\end{lemma}

\begin{proof}
Choose the symplectic lift $I$ of the retained line with
$I^2=-\id_V$.  In the field case use the descended space from
Lemma~\ref{lem:symp-descent}; exact commutation ensures that $I$
is defined over $\F_u$.  It suffices to count, on every descended
plane, determinant-one partners $J_i$ satisfying
$J_i^2=-\id$ and $IJ_i=-J_iI$.

If $-1$ is a square in $\F_u$, choose $r\in\F_u$ with $r^2=-1$
and write $I=\diag(r,-r)$.  Its anticommuting partners are precisely
\[
J_i=\begin{pmatrix}0&a\\-a^{-1}&0\end{pmatrix},
\qquad a\in\F_u^\times,
\]
giving $u-1$ choices.  If $-1$ is nonsquare, identify the plane with
$\F_{u^2}$ so that $I$ is multiplication by a square root of $-1$.
Writing $\tau$ for the nontrivial field automorphism, every
anticommuting map has the form $x\mapsto a\tau(x)$.  Its square is
multiplication by $a\tau(a)$, so the required condition is
$N(a)=-1$.  There are $u+1$ choices.  Their determinant is
$-N(a)=1$, as required.

The choices on the $n$ planes are independent and give
$(u-\chi_u(-1))^n$ global symplectic lifts $J$.  For a fixed $I$
and a resulting projective four-group, the possible partners are
exactly
\[
J,\quad -J,\quad IJ,\quad -IJ.
\]
Indeed the new projective element is one of the two lines of $T$
different from $\langle\overline I\rangle$, and each has its two
symplectic lifts.  Dividing by four gives the stated number of
distinct completions in this family.  In the field case all these
operators were chosen over the descended field, so they commute
exactly with $B$ and remain admissible.  Their projective classes
belong to $L$, so they impose no additional outer-coset condition.
Finally, $u\geq5$ gives $u-\chi_u(-1)\geq4$.

For the second count, fix one admissible $b$ and a representative
$B$ with $B^2=\id_V$.  Every
$D\in\widehat S(\mathcal D)$ commutes exactly with $I,J,B$;
the last assertion follows because $B$ preserves the planes and
fixes their signs.  Thus $(DB)^2=\id_V$, and $DB$ still commutes
exactly with $I,J$.  Its projective class lies in $H\setminus L$
because the class of $D$ lies in $L$.  The $2^n$ linear signs give
$2^{n-1}$ distinct projective axes, with only $D$ and $-D$
identified.  All of them belong to our family, including when their
descended multiplicity forms have different isometry types.
\end{proof}

\subsection{Detection and the strict dimension inequality}

We must check that the parameters counted by the decoration lemma
are coefficients of actual subgroup flags.  This check precedes
the dimension comparison, since different data can have the same
full elementary group $E$.

For any linear group $K=S(\mathcal D)T$, choose symplectic lifts of
its elements.  Their commutators define an alternating bilinear
pairing
\[
\beta_K:K\times K\longrightarrow\{1,-1\},
\qquad [\widehat x,\widehat y]=\beta_K(x,y)\id_V.
\]
The pairing is independent of the lifts.  It vanishes on
$S(\mathcal D)$ and is nondegenerate on $T$ by
\eqref{eq:symp-quaternion}; hence its radical is exactly
$S(\mathcal D)$.  Thus $K$ intrinsically recovers $S(\mathcal D)$,
whose full symplectic preimage has joint eigenspaces precisely the
planes of $\mathcal D$.  Without a field axis, $E=K$; with one,
$K=E\cap M$ is recovered first.  In either case the full group
recovers its plane frame and sign group.

Write $\mathcal L(T)$ for the three lines of $T$.  Its local
coefficient space is
\begin{equation}\label{eq:symp-local-two}
\left\{(c_t)_{t\in\mathcal L(T)}\in\Q^3:
                     \sum_{t\in\mathcal L(T)}c_t=0\right\}.
\end{equation}
It has dimension two and is represented by the coned chains
$\sum_t c_t[t<T]$.  The sum-zero condition cancels the face $[T]$;
dropping the top $T$ leaves one coefficient at each retained line.
This is the quaternion input to Lemma~\ref{lem:quaternion-decoration}.

Let $S_1<\cdots<S_{n-1}=S(\mathcal D)$ be a full sign flag.
Without a field axis, test the split construction on flags
\begin{equation}\label{eq:symp-detection-no-field}
t<T<TS_1<\cdots<TS_{n-1}=E,
\qquad t\in\mathcal L(T).
\end{equation}
A second datum with the same full group has the same sign group
$S$.  Each vertex in its split construction is a product of a
subgroup of $S$ and a subgroup of its chosen complement $T'$.
Since the rank-two vertex $T$ in
\eqref{eq:symp-detection-no-field} is disjoint from $S$, it can
occur only if $T'=T$.  The flags below $T$ then evaluate the three
coordinates in \eqref{eq:symp-local-two}, and intersecting the
remaining vertices with $S$ recovers the sign flag.  The coefficient
on the displayed flag is, up to its fixed shuffle sign, the product
of those two coefficient evaluations.  Such flags therefore detect
the full tensor product of the local and frame parameter spaces.

With a field axis, write $B_0=\langle b\rangle$ and instead test
\begin{equation}\label{eq:symp-detection-field}
B_0<B_0t<B_0T<B_0TS_1<\cdots<B_0TS_{n-1}=E.
\end{equation}
The first vertex is outside $M$.  In a competing split construction,
the only rank-one vertex outside $M$ is its selected field axis, so
that axis must equal $B_0$.  Next,
$(B_0T)\cap M=T$, which is disjoint from $S$ and has rank two.
The same argument as before forces its quaternion complement to
equal $T$.  The choices of $t$ and of the sign flag detect all the
remaining parameters.  Consequently there is no loss of dimension
from coincident cone groups in either construction.

We can now apply the quantitative lemmas.  Let $N$ be the number of
full data and set $a_n=(n-1)!$.  By
Theorem~\ref{thm:frame-bound}, the initial space of actual
coefficients, summed over the fixed decorations, has dimension at
least
\[
2F a_n N,
\qquad
F\geq (17/25)^n,
\qquad F\geq3/8\text{ when }n=3.
\]
Here the same $F$ is a valid lower bound for each multiplicity-space
isometry type.  There are three retained-line incidences per full
datum.  Lemma~\ref{lem:symp-completions} bounds the number of such
partial data by $3N/k_T$, and each costs at most $a_n$ equations.
In the field case the number of data with the field axis removed is
at most $N/k_b$, each costing at most $2a_n$ equations.  Thus
Lemma~\ref{lem:quaternion-decoration} yields a nonzero cycle provided
\begin{equation}\label{eq:symp-inequality}
F>\frac{3}{2k_T}
       +\begin{cases}0,&m=n+1,\\ 2^{1-n},&m=n+2.\end{cases}
\end{equation}
The factor two in the denominator of the quaternion term comes
from the two local parameters in \eqref{eq:symp-local-two}; each
retained-line face has only one.

For $n=3$ the worst completion bound occurs at $u=5$, giving
$k_T=4^3/4=16$.  Even with a field axis,
\begin{equation}\label{eq:symp-smallest}
F\geq\frac38=\frac{12}{32}
   >\frac3{32}+\frac14=\frac{11}{32}.
\end{equation}
For $n=4$ the weaker uniform bound already gives
\[
\left(\frac{17}{25}\right)^4
 =\frac{83521}{390625}
 >\frac{19}{128}
 =\frac{3}{2\cdot4^3}+\frac18.
\]
When $n$ increases by one, the first term on the right is multiplied
by $1/4$ and the second by $1/2$, whereas the left side is
multiplied by $17/25>1/2$.  Induction proves
\eqref{eq:symp-inequality} for every $n\geq4$, and omitting the
field term only strengthens it.

We have therefore obtained a nonzero cycle
\[
\alpha\in\widetilde C_{m-1}(\mathcal A_2(H);\Q)
\]
supported on full flags ending at the rank-$m$ groups $E$ in our
data family.  All its vertices lie in $H$: plane signs and
quaternion generators lie in $L$, and every field axis was chosen
in $H$.  No restriction from a larger automorphism group is needed
in this case.

\subsection{Propagation from a supported full group}

Choose a full flag with nonzero coefficient in $\alpha$, and let
$A'$ be its last group.  The construction gives
$A'\cap L\supseteq A_0(\mathcal D',T')\ne1$ for its datum, and
\[
LA'=H.
\]
Indeed this is immediate if $H=L$; otherwise $A'$ contains its
chosen field involution in the nontrivial coset of $H/L$.  Thus
$LA'$ acts faithfully on $L$, and
$O_2(C_G(LA'))=O_2(C_G(H))=1$.  The group $A'$ belongs to the
family \eqref{eq:symp-configurations} and has its maximum rank $m$.

The selected flag is saturated and has nonzero coefficient in the
cycle $\alpha\in\widetilde C_{m-1}(\mathcal A_2(LA');\Q)$.
Corollary~\ref{cor:maximal-family} therefore applies to $A'$ and the
family \eqref{eq:symp-configurations}.  It contradicts the assumed
minimal counterexample, proving Proposition~\ref{prop:symplectic}.
